% Part: turing-machines
% Chapter: machines-computations
% Section: unary-numbers

\documentclass[../../../include/open-logic-section]{subfiles}

\begin{document}

\olfileid{tur}{mac}{una}
\olsection{સંખ્યાઓનું એક-પ્રતીકી નિરૂપણ}

\begin{explain}
ટ્યુરિંગ મશીનો પોતાની ટેપ પર લખાયેલા પ્રતીકોના અનુક્રમો
પર કામ કરે છે. ટ્યુરિંગ મશીન કઈ વર્ણમાળા વાપરે છે તેના
આધારે આ અનુક્રમો જુદા જુદા ઇનપુટ અને આઉટપુટ દર્શાવી
શકે છે. ખાસ રસ, અલબત્ત, \emph{અંકગણિતીય} વિધેયો,
એટલે પ્રાકૃતિક સંખ્યાઓનાં વિધેયો ગણતા ટ્યુરિંગ મશીનોમાં
છે. ધન પૂર્ણાંકો દર્શાવવાની સરળ રીત તેમને એક જ
પ્રતીક~$\TMstroke$ના અનુક્રમ તરીકે સંકેતિત કરવાની છે.
$n \in \Nat$ હોય, તો $n = 0$ માટે $\TMstroke^n$
ખાલી અનુક્રમ અને અન્યથા બરાબર $n$ $\TMstroke$ ધરાવતો
અનુક્રમ છે.
\end{explain}

\begin{defn}[સંગણના]
ટ્યુરિંગ મશીન~$M$ વિધેય $f\colon \Nat^k \to \Nat$
\emph{ગણે છે} ત્યારે અને ત્યારે જ જ્યારે $M$ ઇનપુટ
\[
\TMstroke^{n_1} \TMblank \TMstroke^{n_2} \TMblank \dots \TMblank \TMstroke^{n_k}
\]
પર થંભે અને આઉટપુટ $\TMstroke^{f(n_1, \dots, n_k)}$
આપે.
\end{defn}

\begin{prob}
  ટ્યુરિંગ મશીન~$M$ વિધેય $f\colon \Nat^k \to \Nat^m$
  ક્યારે ગણે છે તેની વ્યાખ્યા આપો.
\end{prob}

\begin{ex}\ollabel{ex:adder}
\emph{સરવાળો:}
ચાલો $f(n,m) = n + m$ વિધેય ગણતું મશીન બનાવીએ.
તે માટે ટેપ પર $n$ અને~$m$ લંબાઈના $\TMstroke$ના
બે સમૂહો લઈને શરૂ થતું અને $n+m$~$\TMstroke$ના એક
સમૂહ સાથે થંભતું મશીન જોઈએ. ઇનપુટના બે $\TMstroke$
સમૂહોને એક~$\TMblank$ અલગ કરે છે. તેથી એક રીત
$\TMblank$ ધરાવતા ખાનામાં સ્ટ્રોક લખી અને છેલ્લો
$\TMstroke$ ભૂંસી નાખવાની છે.
\begin{figure}\[
\begin{tikzpicture}[->,>=stealth',shorten >=1pt,auto,node distance=2.8cm,
                    semithick]
  \tikzstyle{every state}=[fill=none,draw=black,text=black]

  \node[initial,state]         (A)              {$q_0$};
  \node[state]         (B) [right of=A] {$q_1$};
  \node[state]         (C) [right of=B] {$q_2$};

  \path (A) edge node {\TMtrans{\TMblank}{\TMstroke}{\TMstay}} (B)
           edge [loop above] node {\TMtrans{\TMstroke}{\TMstroke}{\TMright}} (B)
        (B) edge [loop above] node {\TMtrans{\TMstroke}{\TMstroke}{\TMright}} (B)
            edge node {\TMtrans{\TMblank}{\TMblank}{\TMleft}} (C)
        (C) edge [loop above] node {\TMtrans{\TMstroke}{\TMblank}{\TMstay}} (C);
\end{tikzpicture}
\]\caption{$f(x,y) = x+y$ ગણતું મશીન}
\ollabel{fig:adder}
\end{figure}
\end{ex}

\begin{prob}
ઇનપુટ~$\tuple{3,2}$ માટે
\olref[tur][mac][una]{ex:adder}માંના મશીનનાં સંરૂપણોનો
ક્રમ અનુસરો. મશીન $0+0$ ગણે ત્યારે શું થાય છે?
\end{prob}

\begin{explain}
\olref[rep]{ex:doubler}માં આપણે એવું ટ્યુરિંગ મશીન
જોયું હતું જે ઇનપુટ તરીકે $\TMstroke$નો અનુક્રમ લે છે
અને ટેપ પર તેનાથી બમણા $\TMstroke$ના અનુક્રમ સાથે
થંભે છે---બમણું કરનાર મશીન. પરંતુ બમણા કરેલા
$\TMstroke$ના સમૂહની ડાબે આઉટપુટમાં $\TMblank$ રહે
છે. તેથી, તમે કદાચ ધાર્યું હોય તેમ, તે મશીન ખરેખર
$f(x) = 2x$ વિધેય ગણતું નથી. આ ખામી દૂર કરવાની
બે રીતો આપણે વર્ણવીશું.
\end{explain}

\begin{ex}
\olref{fig:doubler-disc}માંનું મશીન $f(x) = 2x$ વિધેય
ગણે છે. \olref[rep]{ex:doubler}માંનું મશીન ઇનપુટ
ભૂંસીને તેના દરેક $\TMstroke$ માટે છેક જમણે બે
$\TMstroke$ લખે છે; તેના બદલે આ મશીન ઇનપુટના
દરેક $\TMstroke$ માટે જમણે એક જ $\TMstroke$ ઉમેરે
છે. ઇનપુટ ક્યાં પૂરું થાય છે તેનો હિસાબ તેને રાખવો
પડે છે. તેથી તે ઇનપુટ અને ઉમેરેલા સ્ટ્રોક વચ્ચે એક
$\TMblank$ રાખે છે, જેને છેવટે $\TMstroke$થી ભરે
છે. આપણે ઇનપુટમાં ક્યાં છીએ તે પણ ``યાદ'' રાખવું
પડે છે, એટલે ઇનપુટના સમૂહમાંના એક $\TMstroke$ને
કામચલાઉ રીતે $\TMblank$થી બદલીએ છીએ.
  \begin{figure}
\[
\begin{tikzpicture}[->,>=stealth',shorten >=1pt,auto,node distance=2.8cm,
                    semithick]
  \tikzstyle{every state}=[fill=none,draw=black,text=black]

  \node[initial,state] (0)              {$q_0$};
  \node[state]         (1) [above of=0] {$q_1$};
  \node[state]         (2) [above of=1] {$q_2$};
  \node[state]         (3) [right of=2] {$q_3$};
  \node[state]         (4) [below of=3] {$q_4$};
  \node[state]         (5) [below of=4] {$q_5$};
  \node[state]         (6) [above right of=4] {$q_6$};
  \node[state]         (7) [below of=6] {$q_7$};
  \node[state]         (8) [below of=7] {$q_8$};

  \path (0) edge node {\TMtrans{\TMstroke}{\TMblank}{\TMright}} (1)
%    (0) edge node {\TMtrans{\TMblank}{\TMblank}{\TMstay}} (h)
    (1) edge [loop left] node {\TMtrans{\TMstroke}{\TMstroke}{\TMright}} (1)
      edge node {\TMtrans{\TMblank}{\TMblank}{\TMright}} (2)
    (2) edge [loop above] node {\TMtrans{\TMstroke}{\TMstroke}{\TMright}} (2)
        edge node {\TMtrans{\TMblank}{\TMstroke}{\TMleft}} (3)
    (3) edge [loop above] node {\TMtrans{\TMstroke}{\TMstroke}{\TMleft}} (3)
        edge node[left] {\TMtrans{\TMblank}{\TMblank}{\TMleft}} (4)
    (4) edge        node[left] {\TMtrans{\TMstroke}{\TMstroke}{\TMleft}} (5)
    (5) edge [loop below] node {\TMtrans{\TMstroke}{\TMstroke}{\TMleft}} (5)
        edge              node {\TMtrans{\TMblank}{\TMstroke}{\TMright}} (0)
    (4) edge      node[sloped] {\TMtrans{\TMblank}{\TMstroke}{\TMright}} (6)
    (6) edge              node {\TMtrans{\TMblank}{\TMstroke}{\TMright}} (7)
    (7) edge [loop right] node {\TMtrans{\TMstroke}{\TMstroke}{\TMright}} (7)
        edge node {\TMtrans{\TMblank}{\TMblank}{\TMleft}} (8)
    (8) edge [loop right] node {\TMtrans{\TMstroke}{\TMblank}{\TMstay}} (8);
    \end{tikzpicture}
\]
\caption{$f(x) = 2x$ ગણતું મશીન}
\ollabel{fig:doubler-disc}
\end{figure}
\end{ex}

\begin{ex}\ollabel{ex:mover}
$f(x) = 2x$ ગણવાની બીજી રીત મૂળ બમણું કરનાર મશીન
રાખીને છેવટે બમણા સ્ટ્રોકના સમૂહને ટેપની છેક ડાબે
લઈ જતી અવસ્થાઓ અને સૂચનાઓ ઉમેરવાની છે.
\olref{fig:mover}માંનું મશીન આ છેલ્લો ભાગ જ કરે છે:
$\TMblank$ના સમૂહ પછી $\TMstroke$નો સમૂહ હોય
એવી ટેપ પર (અને હેડ $\TMblank$ના સમૂહમાં ગમે ત્યાં
હોય) તેને શરૂ કરીએ, તો તે એક પછી એક $\TMstroke$
ભૂંસીને ટેપની શરૂઆતમાં લખે છે. કામ પૂરું થયું છે તે
ઓળખી શકાય એ માટે તે પહેલાં $\TMstroke$ના સમૂહનો
છેડો $\TMendtape$ પ્રતીક વડે ચિહ્નિત કરે છે; છેવટે
આ પ્રતીક ભૂંસાઈ જાય છે. આપણે અવસ્થાઓને~$q_6$થી
ક્રમ આપવાનું શરૂ કર્યું છે, જેથી તેને બમણું કરનાર મશીનમાં
ઉમેરી શકાય. તમારે વધારાની સૂચના
$\delta(q_0, \TMblank) = \tuple{q_6,\TMblank,\TMstay}$,
એટલે~$q_0$થી~$q_6$ સુધીનું
$\TMtrans{\TMblank}{\TMblank}{\TMstay}$ ચિહ્નવાળું તીર,
ઉમેરવાનું છે. (એક સૂક્ષ્મ સમસ્યા છે: મળેલું મશીન
ઇનપુટ~$x=0$ માટે કામ કરતું નથી. તેને આપણે અભ્યાસપ્રશ્ન
તરીકે રાખીએ છીએ.)
\begin{figure}
  \[
  \begin{tikzpicture}[->,>=stealth',shorten >=1pt,auto,node distance=2.8cm,
                      semithick]
    \tikzstyle{every state}=[fill=none,draw=black,text=black]
    \node[initial,state] (6)              {$q_6$};
    \node[state]         (7) [right of=6] {$q_7$};
    \node[state]         (8) [right of=7] {$q_8$};
    \node[state]         (9) [below of=8] {$q_9$};
    \node[state]         (10) [left of=9] {$q_{10}$};
    \node[state]         (11) [left of=10]  {$q_{11}$};
    \node[state]         (12) [below of=11]  {$q_{12}$};
    \node[state]         (13) [right of=12] {$q_{13}$};
    \node[state]         (14) [right of=13] {$q_{14}$};
    \path
    (6)  edge node {\TMtrans{\TMblank}{\TMblank}{\TMright}} (7)
    (7)  edge [loop above] node {\TMtrans{\TMblank}{\TMblank}{\TMright}} (7)
         edge node {\TMtrans{\TMstroke}{\TMstroke}{\TMright}} (8)
    (8)  edge [loop above] node {\TMtrans{\TMstroke}{\TMstroke}{\TMright}} (8)
         edge node {\TMtrans{\TMblank}{\TMendtape}{\TMleft}} (9)
    (9)  edge [loop right] node {\TMtrans{\TMstroke}{\TMstroke}{\TMleft}} (9)
         edge node {\TMtrans{\TMblank}{\TMblank}{\TMright}} (10)
    (10) edge node[above] {\TMtrans{\TMstroke}{\TMblank}{\TMleft}} (11)
    (11) edge [loop above] node {\TMtrans{\TMblank}{\TMblank}{\TMleft}} (11)
         edge node[left] {\begin{tabular}{@{}l@{}}
          \TMtrans{\TMendtape}{\TMendtape}{\TMright}\\
          \TMtrans{\TMstroke}{\TMstroke}{\TMright}
         \end{tabular}} (12)
    (12) edge node[] {\TMtrans{\TMblank}{\TMstroke}{\TMright}} (13)
    (13) edge [loop below] node {\TMtrans{\TMblank}{\TMblank}{\TMright}} (13)
         edge node[sloped] {\TMtrans{\TMstroke}{\TMblank}{\TMleft}} (11)
    (13) edge node {\TMtrans{\TMendtape}{\TMblank}{\TMstay}} (14);
  \end{tikzpicture}
  \]
  \caption{$\TMstroke$ના સમૂહને ડાબે ખસેડવું}
  \ollabel{fig:mover}
\end{figure}
\end{ex}

\begin{prob}
\olref[tur][mac][una]{ex:mover}માં આપણે
\olref[rep]{ex:doubler}માંના બમણું કરનાર મશીન અને
\olref[tur][mac][una]{fig:mover}માંના ખસેડનાર મશીનને
જોડીને બનેલા મશીનનું વર્ણન કર્યું. આ સંયુક્ત મશીનને
ઇનપુટ~$x=0$, એટલે ખાલી ટેપ, પર શરૂ કરીએ તો શું
થાય છે? આ કિસ્સામાં મશીન આઉટપુટ~$2x=0$ સાથે
થંભે એ માટે તેને કેવી રીતે સુધારશો? (એક અવસ્થા અને
એક સંક્રમણ ઉમેરીને આ કરી શકવું જોઈએ.)
\end{prob}

\sourcecorrection{OLTM-003}{સ્થિર અંગ્રેજી મૂળના આ પ્રશ્નમાં
બમણું કરનાર મશીન માટે $f(x)=2x$ આપતી
\olref[tur][mac][una]{fig:doubler-disc} આકૃતિનો સંદર્ભ છે.
પરંતુ તરત અગાઉનું વર્ણન મૂળ
\olref[rep]{ex:doubler} મશીનને
\olref[tur][mac][una]{fig:mover} સાથે જોડે છે; પ્રશ્નમાં
તે જ મૂળ મશીનનો સંદર્ભ આપ્યો છે.}

\begin{prob}
\emph{બાદબાકી:} $n$ અને~$m$ લંબાઈની બે ખાલી ન હોય
એવી સ્ટ્રોકની શબ્દમાળાઓ ઇનપુટ તરીકે મળે, જ્યાં
$n > m$ હોય, ત્યારે $f(n,m) = n - m$ વિધેય ગણતું
ટ્યુરિંગ મશીન રચો.
\end{prob}

\begin{prob}
\emph{સમાનતા:} નીચેનું વિધેય ગણતું ટ્યુરિંગ મશીન રચો:
\[
\fn{equality}(n,m) =
\begin{cases}
  \text{1} & \text{જો~$n = m$ હોય} \\
  \text{0} & \text{જો~$n \neq m$ હોય}
\end{cases}
\]
અહીં~$n$ અને~$m \in \PosInt$ છે.
\end{prob}

\begin{prob}
$x$ અને $y$ ધન પૂર્ણાંકોને ટેપ પર $\TMblank$થી
અલગ કરેલી $\TMstroke$ની શબ્દમાળાઓ વડે દર્શાવ્યા હોય,
ત્યારે $\min(x,y)$ વિધેય ગણતું ટ્યુરિંગ મશીન રચો.
મશીનની વર્ણમાળામાં વધારાનાં પ્રતીકો વાપરી શકો છો.

વિધેય $\min$ તેના આર્ગ્યુમેન્ટોમાંથી સૌથી નાનું મૂલ્ય
પસંદ કરે છે: તેથી $\min(3,5)=3$, $\min(20,16)=16$
અને $\min(4,4)=4$ વગેરે.
\end{prob}

\begin{defn}
  ટ્યુરિંગ મશીન~$M$ આંશિક વિધેય $f\colon \Nat^k
  \pto \Nat$ ગણે છે ત્યારે અને ત્યારે જ જ્યારે
  \begin{enumerate}
    \item $f(n_1, \dots, n_k) = m$ હોય તો $M$ ઇનપુટ
    $\TMstroke^{n_1}\concat\TMblank\concat
    \dots \concat\TMblank\concat\TMstroke^{n_k}$ પર
    આઉટપુટ $\TMstroke^{m}$ સાથે થંભે છે.
    \item $f(n_1, \dots, n_k)$ અવ્યાખ્યાયિત હોય તો
    $M$ થંભતું જ નથી, અથવા $\TMstroke$ના એક જ
    સમૂહથી ન બનેલા આઉટપુટ સાથે થંભે છે.
  \end{enumerate}
\end{defn}

\end{document}
